\pagina{Contraste con resultados publicados del FHCP}
El contraste bibliográfico reutiliza el inventario y los localizadores verificados de M6. Las fuentes primarias no proporcionan condiciones equivalentes a nuestras ejecuciones: difieren métodos, instancias, presupuestos, equipos y repeticiones. Siguiendo buenas prácticas de comparación de optimización, esas diferencias se explicitan antes de interpretar cifras \citep{beiranvand2017BestPractices}.

\subsection{Cobertura del concurso y dificultad estructural}
El portal oficial informa que Cohen y Coudert presentaron ciclos para 985 de los 1001 grafos del concurso \citep{flindersFhcpChallenge}. Ese resultado es cobertura acumulada durante una ventana de concurso; no equivale a diez o treinta semillas con un plazo individual homogéneo. Comparar su porcentaje directamente con 15 de 149 grafos iniciados produciría una equivalencia que los datos no sostienen.

La construcción del conjunto consideró dificultad para Concorde, LKH y SLH, seleccionando candidatos problemáticos para al menos dos de esos métodos \citep{haythorpe2018Fhcp}. Esta procedencia ayuda a interpretar por qué un constructor o movimiento local puede estancarse, pero no aporta por sí sola una matriz temporal reproducible para esos tres solvers. Los tours oficiales sirven como referencias de soluciones conocidas y no como mediciones propias de rendimiento.

\subsection{Contraste por instancia y tipo de salida}
Cohen y Coudert describen un oráculo de programación entera con SageMath y CPLEX y publican tiempos puntuales para algunos IDs; graph1 y graph2 aparecen con 0.76 y 0.01 segundos \citep{cohenCoudert2017Challenge}. No se precisan todas las condiciones necesarias para una comparación temporal equivalente. En nuestra campaña histórica, los IDs compartidos pertenecían a comprobación o verificación, no a evaluación final; el CSV trazable conserva esa distinción en 36 contrastes.

HybridHAM combina construcción voraz, transformaciones por rotación y criterios de alcanzabilidad \citep{seeja2018HybridHam}. Su tabla diferencia HP y HC: graph1 aparece como camino con 0.0625 segundos y graph2 como ciclo con 0.0469. El artículo indica Core i5, 4 GB de RAM y MATLAB 13RA, sin especificar todas las condiciones de repetición y presupuesto requeridas para comparar con nuestro entorno.

Además, su texto de la sección 4.3 declara seis ciclos entre 250 instancias y 75 caminos, mientras la tabla 4 etiqueta más de seis instancias como HC. Se conserva esa discrepancia interna y no se infiere un total corregido ni una tasa global. Los tiempos por ID se atribuyen a las filas correspondientes; no se convierte un HP en ciclo para aumentar el éxito publicado.

